%% kapa: exam=1 round=31 year=2020 period=1 form=A subject=economics subject_ko=경제학원론 seq=30 number=70
%% source: Q-Net 감정평가사 기출문제 (cst00309) artlSeq=5207219 file=extracted/1008/5207219/A형/A형/2020년도 제31회 감정평가사 1차 시험_1교시_A형.pdf
%% conversion: kordoc 4.24.51 -> tools/gen1.py (auto); review=proofread-sonnet; flags=none
\begin{question}{70}
총수요 충격 및 총공급 충격에 관한 설명으로 옳지 않은 것은? (단, 총수요곡선은 우하향, 총공급곡선은 우상향)
\begin{choices}
  \item 총수요 충격으로 인한 경기변동에서 물가는 경기순행적이다.
  \item 총공급 충격으로 인한 경기변동에서 물가는 경기역행적이다.
  \item 총공급 충격에 의한 스태그플레이션은 합리적 기대 가설이 주장하는 정책무력성의 근거가 될 수 있다.
  \item 명목임금이 하방 경직적일 경우 음(-)의 총공급 충격이 발생하면 거시경제의 불균형이 지속될 수 있다.
  \item 기술진보로 인한 양(+)의 총공급 충격은 자연실업률 수준을 하락시킬 수 있다.
\end{choices}
\end{question}
\answer{3}
